\clearpage
\chapter*{Erratum à l'exposé n\textsuperscript{o} 190}
\addcontentsline{toc}{chapter}{Erratum à l'exposé n\textsuperscript{o} 190}
\setcounter{footnote}{0}
\src{300}

\noindent\textbf{3. Technique de descente et théorèmes d'existence en
géométrie algébrique [TDTE I]. I : Généralités. Descente par morphismes
fidèlement plats.}

\begin{center}
[Séminaire Bourbaki, t. 12, 1959/60, n\textsuperscript{o} 190, 29 p.]
\end{center}

\noindent\textbf{Page 1, ligne 10.} --- Il semble maintenant excessif de dire
que la technique de descente est « à la base de la plupart des théorèmes
d'existence en géométrie algébrique ». Cela est vrai dans une large mesure
pour les techniques non projectives faisant l'objet des deux premiers exposés
de la série TDTE I à VI, mais non pour les techniques projectives (TDTE IV, V,
VI).

\noindent\textbf{Page 5, ligne 16.} --- Il est inutile de supposer que $\alpha$
soit un morphisme de $\mathcal F$-descente.

\noindent\textbf{Page 20, remarque.} --- Un morphisme
$S'\longrightarrow S$ quasi-fini, étale, surjectif, ou un morphisme
$\operatorname{Spec}(\widehat A)\longrightarrow\operatorname{Spec}(A)$ n'est
pas toujours un morphisme de descente strict, même si $A$ est un anneau local
d'une courbe algébrique sur un corps algébriquement clos $k$, et
$S=\operatorname{Spec}(A)$. Ainsi on peut trouver un morphisme propre et simple
$f:X\longrightarrow S$, faisant de $X$ un fibré principal de base $S$ sous une
courbe elliptique $E$, tel que $f':X'\longrightarrow S'$ soit projectif, mais
$f$ n'étant pas projectif. C'est donc en même temps un exemple d'un fibré
principal homogène non isotrivial sous un schéma abélien.

\noindent\textbf{Page 26, ligne 9.} --- Lire « Chow--Lang » au lieu de
« Igusa--Lang ».

Pour divers détails touchant la théorie de la descente, voir SGA VI, VII,
VIII.

\par\smallskip
\begingroup\footnotesize
\endgroup
